problem:  
Let \( M^8 \) be a compact, simply connected Riemannian manifold with **nonnegative sectional curvature** admitting an **effective, isometric \( S^1 \)-action**.  
Let \( Q = M / S^1 \) denote its **7‑dimensional Alexandrov space** of curvature ≥ 0, and assume \( Q \) is **homeomorphic to \( S^7 \)**.  

**Assume further** that the image \( \Sigma \subset Q \) of the nonprincipal orbits forms a **connected, codimension‑2 extremal subset** that is **homeomorphic to \( S^5 \)** (this is a conditional structural hypothesis, not a general fact about all such actions).

**Question (Conditional Curvature–Knot Rigidity Problem in Dimension (8, 7)):**  
Under the above assumptions, can \( \Sigma \subset Q \) be a **topologically knotted embedding** \( S^5 \hookrightarrow S^7 \)?  
Equivalently, is it possible for such a quotient Alexandrov \( S^7 \) of nonnegative curvature, arising from an isometric circle action on a simply connected, nonnegatively curved \( 8 \)-manifold, to contain a **knotted codimension‑2 extremal subset** homeomorphic to \( S^5 \)?  
Or must every such subset be **unknotted**—that is, ambiently isotopic to the standard equatorial \( S^5 \subset S^7 \) appearing in linear models of circle actions on \( S^8 \)?  

Why is it a "good" problem:  
This formulation corrects earlier logical ambiguities by framing the structure of the singular set as a **conditional hypothesis** rather than an automatic property, ensuring mathematical validity. It remains a deep, open question at the confluence of geometry, topology, and symmetry:

* **Simplicity with depth:**  The statement is short—asking whether curvature ≥ 0 enforces unknottedness of an \( S^5 \) inside an Alexandrov \( S^7 \)—yet encodes intricate geometric and topological constraints.  
* **Interdisciplinary bridge:**  
  - From **Riemannian geometry** and **Alexandrov theory**, it probes whether curvature restrictions can enforce global topological simplicity on singular orbits.  
  - From **transformation group theory**, it challenges the extent to which circle actions on nonnegatively curved manifolds must be equivariantly equivalent to linear ones.  
  - From **high‑dimensional knot theory**, it connects to the classification of smooth embeddings \( S^5 \hookrightarrow S^7 \), whose isotopy types are rich and related to the framed cobordism group \( \pi_5(SO(2)) \simeq 0 \) and higher homotopy features of \(S^2\)-bundles, though no known structure theorem links them to curvature geometry.  
* **Open and foundational:**  Nothing is currently known about whether an Alexandrov space of curvature ≥ 0 can realize a *knotted* codimension‑2 extremal subset under such a symmetry constraint.  
* **Potential significance:**  
  - A positive “unknottedness” result would demonstrate a new **curvature‑driven rigidity**—that nonnegative curvature and circle symmetry eliminate topological complexities in orbit strata.  
  - A counterexample would provide the first instance of a **knotted extremal subset** compatible with curvature ≥ 0, revolutionizing our understanding of curvature’s topological flexibility.  

By combining analytical curvature bounds, global group‑action geometry, and embedding topology into one concise yet profound question, this problem exemplifies the essence of a “good” research problem—elegant, cross‑cutting, and capable of advancing multiple mathematical frontiers simultaneously.